% Templat TA TIF UIN Suska Riau -- (c) 2026 Benny Sukma Negara
% Bab ini hanya dicetak pada opsi [laporan].
\chapter{Pembahasan}
\label{bab:pembahasan}

%==========================================================
\section{Deskripsi}
\label{sec:deskripsi}

Bab Pembahasan berisi semua pembahasan yang dibutuhkan terkait Tugas Akhir. Subbab dapat disesuaikan dengan kebutuhan Laporan Tugas Akhir dengan tetap mengedepankan prinsip-prinsip:
\begin{enumerate}
  \item Akurat,
  \item Jelas,
  \item Berurutan.
\end{enumerate}

Adapun subbab yang disarankan (dapat disesuaikan) antara lain:
\begin{enumerate}
  \item Analisa
  \item Perancangan
  \item Implementasi
  \item Pengujian
\end{enumerate}

Adapun hal-hal yang perlu diperhatikan terkait dengan Bab Pembahasan antara lain:
\begin{enumerate}
  \item Tidak perlu menjelaskan tentang teori/konsep/metode/algoritma, cukup dirujuk ke sitasi yang relevan atau ke subbab yang ada di Bab Kajian Pustaka.
\end{enumerate}

Untuk struktur penomoran pada bagian isi laporan dapat menggunakan penomoran bertingkat dengan urutan sebagai berikut:
\begin{enumerate}
  \item Level pertama menggunakan penomoran latin, seperti 1, 2, 3 dan seterusnya.
  \item Level kedua menggunakan penomoran alfabet (huruf kecil), seperti contoh berikut:
  \begin{enumerate}
    \item Ini contoh penomoran level kedua
    \item Ini contoh penomoran level kedua
    \item Ini contoh penomoran level kedua
  \end{enumerate}
  \item Level ketiga menggunakan penomoran romawi (huruf kecil), seperti contoh berikut:
  \begin{enumerate}
    \item Ini contoh penomoran level kedua
    \item Ini contoh penomoran level kedua
    \begin{enumerate}
      \item Ini contoh penomoran level ketiga
      \item Ini contoh penomoran level ketiga
      \item Ini contoh penomoran level ketiga
      \item Ini contoh penomoran level ketiga
    \end{enumerate}
  \end{enumerate}
  \item Maksimal tingkat penomoran hanya sampai level ketiga. Mahasiswa harus membuat sedemikian rupa sehingga tingkat penomoran hanya sampai level ketiga.
\end{enumerate}

Untuk format penyisipan gambar dapat dilihat pada contoh berikut:

\begin{figure}[H]
\centering
\gambar[0.45\textwidth]{logo-uin-suska.png}
\caption{Tampilan aplikasi pengolah dokumen}
\label{fig:contohgambar}
\end{figure}

Sedangkan untuk penyisipan tabel dapat dilihat pada contoh berikut:

\begin{table}[H]
\caption{Contoh penyisipan tabel}
\label{tab:contohtabel}
\begin{tabularx}{\textwidth}{|p{10mm}|X|p{25mm}|p{25mm}|}
\hline
No & Judul           & Status & Jumlah\\ \hline
1  & Teks pada tabel & Baik   & 124\\ \hline
2  & Teks pada tabel & Baik   & 124\\ \hline
3  & Teks pada tabel & Baik   & 124\\ \hline
\end{tabularx}
\end{table}

Sedangkan untuk penyisipan rumus matematika dapat dilihat pada contoh berikut:
\begin{rumus}{Rumus Satu}
  x = \frac{-b \pm \sqrt{b^2 - 4ac}}{2a}
  \label{eq:abc}
\end{rumus}
